% Ad Atticum 15.7 --- headnote
% ad-atticum-15-07, 28 or 29 May 44 BC, in Tusculano

\textit{Cicero to Atticus, written at the Tusculan villa on
either 28 or 29 May 44 BC --- Perseus dateline \textit{Scr. in
Tusculano v aut iv K. Iun. a. 710 (44)}. (The works.yaml entry
takes the later of the two days, 29 May; the dateline itself is
indeterminate.) A very short acknowledgement: a single
paragraph, with no Perseus section divisions. Cicero thanks
Atticus for forwarded letters, singling out one from
``our Sextus'' --- Sextus Pompeius, then in Spain --- whose
sentiments on the commonwealth and whose style of writing
pleased him even before he came to the passage in which he was
praised.}

\textit{The closing sentences turn dismissive on Servius
Sulpicius Rufus, ``our peacemaker,'' who seems to have set off
on his embassy in the company of a little clerk and to be
afraid of every legal quibble. The figure is technical Roman
legal language: \textit{ex iure manum consertum} is the
formula by which two parties to a property dispute laid hands
on the thing in question to begin litigation; Cicero says
Servius ought to have proceeded not by that initial formal
claim but by what comes after --- that is, with substance
rather than ceremony. The Latin around \textit{cum sensus eius
de re publica} carries a daggered crux in the manuscripts; the
sense is preserved.}
